% Pages 41–50

% ---- Pages 41: Problems 2.23–2.25 ----

\begin{enumerate}
  \item[b)] Converting $10^9$ years to seconds:
  \[
    10^9 \,\text{yr} \times \frac{52\,\text{wk}}{\text{yr}} \times \frac{7\,\text{days}}{\text{wk}}
    \times \frac{24\,\text{hr}}{\text{day}} \times \frac{60\,\text{min}}{\text{hr}}
    \times \frac{60\,\text{s}}{\text{min}}
    = 10^9 \times 365.25 \times 86400 \approx 10^{16}\,\text{s}.
  \]

  \item[c)] ``Accessible'' means physically allowed by the conservation laws, not that the system will
  actually visit every such state. In practice, not every accessible microstate will be reached.
\end{enumerate}

\textbf{2.24.}
\begin{enumerate}
  \item[a)] As computed in Problem~2.22, the peak multiplicity for a two-state paramagnet is
  $\Omega\!\left(N,\tfrac{N}{2}\right) = 2^N\!\sqrt{\tfrac{2}{\pi N}} = \Omega_{\max}$.

  \item[b)] For a general spin-up count $N_0$,
  \[
    \Omega(N, N_0) = \binom{N}{N_0} = \frac{N!}{N_0!\,(N-N_0)!}.
  \]
\end{enumerate}

\textbf{2.25.}
\begin{enumerate}
  \item[a)] The random walker ends back at the center. In terms of steps, the sequence is
  $\epsilon \to \epsilon \to \epsilon \to \epsilon \leftarrow \epsilon$, illustrating that a return
  to the origin is the most probable outcome.
\end{enumerate}

% ---- Page 42: Section 2.5 — Ideal Gas ----

\section{The Ideal Gas}

\subsection{Multiplicity of a Monatomic Ideal Gas}

The central lesson so far is that only a tiny fraction of the microstates of a large interacting
system carry appreciable probability; the rest are overwhelmingly unlikely. We now apply this
reasoning to a continuous system: the ideal gas.

Consider first a monatomic ideal gas consisting of a single atom in a fixed volume $V$ with fixed
total energy $U$. The multiplicity of this system has two independent sources: the atom can be
anywhere in the spatial volume $V$, and its momentum vector can point anywhere in the region of
momentum space consistent with the energy constraint. Hence
\[
  \Omega_1 \propto V \cdot V_p,
\]
where $V_p$ denotes the ``volume'' accessible in momentum space. The energy constraint
\[
  U = \frac{1}{2m}(p_x^2 + p_y^2 + p_z^2) \quad\Longrightarrow\quad p_x^2 + p_y^2 + p_z^2 = 2mU
\]
restricts the momentum to the \emph{surface} of a sphere of radius $\sqrt{2mU}$ in
three-dimensional momentum space.

To count microstates properly we need quantum mechanics, which discretises the states and
provides a natural lower bound on the uncertainty of any measurement.

\begin{fact}
\textbf{Heisenberg Uncertainty Principle.} For a particle confined to a region of length $L_x$ in
position space, the uncertainty in position and momentum satisfies $\Delta x \,\Delta p_x \gtrsim h$,
where $h$ is Planck's constant.
\end{fact}

The number of distinct position states in the $x$-direction is $L_x/\Delta x$, and the number of
distinct momentum states is $L_p/\Delta p_x$. The total number of distinct states along $x$ is
therefore $(L_x/\Delta x)(L_p/\Delta p_x) = L_x L_p / h$. Extending this to three dimensions gives
the fundamental result:

\begin{center}
[DIAGRAM: Position space showing a wavefunction of width $\Delta x$ inside a box of length $L_x$;
momentum space showing a wavefunction of width $\Delta p_x$ inside a range $L_p$]
\end{center}

\begin{formula}
The multiplicity of a single monatomic molecule in three dimensions is
\[
  \Omega_1 = \frac{V\, V_p}{h^3},
\]
where $V_p$ is the accessible volume in momentum space.
\end{formula}

\subsection{Multiplicity for $N$ Molecules}

For two molecules the energy constraint becomes
\[
  \|\vb{p}_1\|^2 + \|\vb{p}_2\|^2 = 2mU,
\]
which describes the surface of a six-dimensional hypersphere. Because the molecules move
independently, the spatial part of the multiplicity factorises, and we write
\[
  \Omega_2 = \frac{V^2}{h^6} \cdot A_{\text{mom}},
\]
where $A_{\text{mom}}$ is the area of the six-dimensional momentum hypersphere. This formula
treats the molecules as distinguishable. If they are indistinguishable—as identical molecules
must be—we divide by $2!$ to avoid overcounting:
\[
  \Omega_2 = \frac{1}{2}\cdot\frac{V^2}{h^6}\cdot(\text{area of momentum hypersphere}).
\]

Extending to $N$ indistinguishable molecules, the momentum constraint restricts the $3N$-component
momentum vector to the surface of a $3N$-dimensional hypersphere of radius $\sqrt{2mU}$. The
surface area of a $d$-dimensional sphere of radius $r$ is
\[
  A_d = \frac{2\pi^{d/2}}{\left(\tfrac{d}{2}-1\right)!}\,r^{d-1}.
\]
Applying this with $d = 3N$ and $r = \sqrt{2mU}$:

\begin{formula}
The multiplicity of an ideal gas of $N$ indistinguishable monatomic molecules is
\[
  \Omega_N = \frac{1}{N!}\cdot\frac{V^N}{h^{3N}}\cdot
    \frac{2\pi^{3N/2}}{\!\left(\tfrac{3N}{2}-1\right)!}\,(\sqrt{2mU})^{3N-1}
  \;\approx\; \frac{1}{N!}\cdot\frac{V^N}{h^{3N}}\cdot
    \frac{\pi^{3N/2}}{\left(\tfrac{3N}{2}\right)!}\,(2mU)^{3N/2},
\]
where the second form drops the subleading $r^{-1}$ factor, valid for large $N$.
\end{formula}

Isolating the dependences on $U$ and $V$, the multiplicity takes the compact form
\begin{formula}
\[
  \Omega(U,V,N) = f(N)\cdot V^N \cdot U^{3N/2}.
\]
\end{formula}
The exponent $3N/2$ is the number of quadratic degrees of freedom divided by two; it is not a
universal result but is specific to the monatomic ideal gas.

\subsection{Two Interacting Ideal Gases}

\begin{center}
[DIAGRAM: Two boxes side by side separated by a partition that allows energy (but not matter or
volume) to pass through; left box labelled $U_A,\,V_A,\,N$; right box labelled $U_B,\,V_B,\,N$]
\end{center}

Suppose two ideal gases are separated by a partition that permits only energy exchange. The total
multiplicity is the product $\Omega_{\text{total}} = \Omega_A\,\Omega_B$, which gives
\[
  \Omega_{\text{total}} = \bigl[f(N)\bigr]^2 (V_A V_B)^N (U_A U_B)^{3N/2}.
\]
This has the same mathematical structure as the multiplicity of two interacting Einstein solids,
so all of the earlier analysis carries over directly. By arguments analogous to those in
Section~2.4, the width of the multiplicity peak about the most likely energy sharing is
\[
  \Delta U \sim \frac{U_{\text{total}}}{\sqrt{3N/2}},
\]
confirming that for macroscopic $N$ the peak is extraordinarily sharp.

The same logic applies if instead of energy we allow \emph{volume} to be exchanged through a
movable partition.

\begin{center}
[DIAGRAM: Same two-box setup with a movable partition; the question is which value of $V_A$
maximises $\Omega_{\text{total}}$]
\end{center}

By symmetry, the most likely split is equal volumes, and the width of the peak around that split is
\[
  \Delta V \sim \frac{V_{\text{total}}}{\sqrt{N}}.
\]

\textbf{Example.} What is the probability that all $N$ molecules of a gas spontaneously occupy
only the left half of their container?

\begin{center}
[DIAGRAM: Box divided by a dashed line; all $N$ molecules on the left half; arrow to same box with
molecules distributed throughout]
\end{center}

Confining all molecules to the left half is equivalent to having an ideal gas in a volume $V/2$.
Since $\Omega \propto V^N$, shrinking the volume by half reduces the multiplicity by $2^N$, so the
probability of this configuration is $2^{-N}$—utterly negligible for any macroscopic $N$.

\subsection*{Problem}

If the volume changes from $V$ to $0.99V$, the multiplicity changes by
\[
  \frac{\Omega_{\text{new}}}{\Omega_{\text{orig}}} = (0.99)^N,
\]
which for $N \sim 10^{23}$ is fantastically small.

% ---- Page 45: Section 2.6 — Entropy ----

\section{Entropy}

The pattern that has emerged across every system we have studied is that any large system will
almost certainly be found in—or very near—the macrostate of greatest multiplicity. Fluctuations
away from this peak are not forbidden, only exponentially suppressed. This observation is the
statistical underpinning of the second law of thermodynamics.

\begin{definition}
The \textbf{entropy} of a macrostate with multiplicity $\Omega$ is
\[
  S \equiv k\ln\Omega,
\]
with units of J/K. Because $\Omega \geq 1$, entropy is non-negative. Because $\Omega =
\Omega_A\Omega_B$ for independent subsystems, entropy is additive: $S_{\text{total}} = S_A + S_B$.
\end{definition}

\begin{theorem}
\textbf{Second Law of Thermodynamics.} The total entropy of an isolated system tends to increase
(or remain constant) over time. Equivalently, the multiplicity tends to increase.
\end{theorem}

\textbf{Example: high-temperature Einstein solid.} In the limit $q \gg N$, the multiplicity of
an Einstein solid is $\Omega \approx (eq/N)^N$, so
\[
  S = k\ln\!\left[\left(\frac{eq}{N}\right)^{\!N}\right] = kN\ln\!\left(\frac{eq}{N}\right)
  = kN\!\left(\ln\frac{q}{N} + 1\right).
\]

The additive property of entropy is worth emphasising. When two subsystems $A$ and $B$ are
independent, $\Omega_{\text{total}} = \Omega_A\Omega_B$, and therefore
\[
  S_{\text{total}} = k\ln(\Omega_A\Omega_B) = k\ln\Omega_A + k\ln\Omega_B = S_A + S_B.
\]
This holds when the macrostates of $A$ and $B$ are specified separately; it requires care when $A$
and $B$ interact, because then $\Omega_{\text{total}}$ must be computed directly as a sum over
shared microstates.

\begin{center}
[DIAGRAM: Left: single box containing both gases $A$ and $B$ together. Right: gas $A$ and gas $B$
in separate boxes. The entropy formula $S_A + S_B$ applies to the separated case when macrostates
are independently specified.]
\end{center}

In practice, for a system with an overwhelmingly sharp multiplicity peak, the total entropy is well
approximated by $S \approx k\ln\Omega_{\text{most likely}}$.

% ---- Page 46: Sackur–Tetrode, Free Expansion, Entropy of Mixing ----

\subsection{Entropy of an Ideal Gas: the Sackur--Tetrode Equation}

Starting from the multiplicity
\[
  \Omega_N \approx \frac{1}{N!}\cdot\frac{V^N}{h^{3N}}\cdot
    \frac{\pi^{3N/2}}{\left(\tfrac{3N}{2}\right)!}\,(2mU)^{3N/2}
\]
and applying Stirling's approximation ($\ln n! \approx n\ln n - n$) to both $N!$ and
$\left(\tfrac{3N}{2}\right)!$, one obtains the entropy of an ideal gas:

\begin{formula}
\textbf{Sackur--Tetrode Equation.}
\[
  S(N,V,U) = Nk\!\left[\ln\!\left(\frac{V}{N}\left(\frac{4\pi m U}{3Nh^2}\right)^{3/2}\right)
  + \frac{5}{2}\right].
\]
\end{formula}

As expected, $S$ increases with $N$, $V$, or $U$, and with the energy per particle $U/N$.
For a volume change at fixed $U$ and $N$, the Sackur--Tetrode equation gives
\begin{formula}
\[
  \Delta S = S(V_f) - S(V_i) = Nk\ln\!\frac{V_f}{V_i},
\]
\end{formula}
which applies, for example, in a quasistatic isothermal compression or expansion.

\subsection{Free Expansion}

\begin{center}
[DIAGRAM: Gas confined to the left half of an insulated box by a removable partition; arrow
pointing right; gas now fills the entire box]
\end{center}

In a free expansion, a partition is removed and the gas expands into vacuum. No work is done and
no heat is transferred, so $U$ is unchanged. Yet the volume doubles, and by the Sackur--Tetrode
equation the entropy increases by $Nk\ln 2 > 0$. This is an irreversible process: entropy has been
created from nothing. Unlike a quasistatic expansion, there is no external reservoir to absorb
that entropy, and the process cannot be run in reverse without violating the second law.

\subsection{Entropy of Mixing}

Allowing two \emph{different} ideal gases to mix also generates entropy. Suppose gases $A$ and
$B$ each contain $N$ molecules, occupy volume $V$, and have the same internal energy. When the
partition between them is removed, each gas expands from $V$ to $2V$, so
\[
  \Delta S_A = Nk\ln\frac{2V}{V} = Nk\ln 2 = \Delta S_B.
\]

\begin{center}
[DIAGRAM: Two separate boxes (gas $A$ left, gas $B$ right) with arrow pointing to a single box
containing both gases mixed]
\end{center}

\begin{formula}
The \textbf{entropy of mixing} for two distinct ideal gases is
\[
  \Delta S_{\text{mix}} = \Delta S_A + \Delta S_B = 2Nk\ln 2.
\]
\end{formula}

This result applies only when the gases are \emph{different}. Doubling a gas of a single species
does not produce entropy of mixing; the entropy of $2N$ molecules of helium in volume $2V$ is
simply $2S_{\text{initial}}$, not $2S_{\text{initial}} + 2Nk\ln 2$.

\begin{center}
[DIAGRAM: Two boxes each containing gas $A$ (same species); arrow with $\neq$ sign; single larger
box of gas $A$, illustrating no entropy of mixing for identical species]
\end{center}

The reason lies in the $1/N!$ factor in the multiplicity. The molecules of a single-species gas
are \emph{indistinguishable}: swapping two helium atoms produces no new microstate. The $1/N!$
accounts for this, so combining two identical volumes does not open up any new microstates beyond
what doubling $N$ already accounts for. For two \emph{different} species, swapping an $A$
molecule with a $B$ molecule genuinely creates a new microstate, and the $1/N!$ factor within
each species does not eliminate that additional count.

\begin{center}
[DIAGRAM: Two identical molecules in a box with a double-headed arrow, indicating that swapping
them yields the same microstate, justifying the $1/N!$ factor]
\end{center}

\subsection{Reversible and Irreversible Processes}

\begin{definition}
An \textbf{irreversible} process increases the total entropy of the universe; a \textbf{reversible}
process leaves the total entropy unchanged.
\end{definition}

Free expansion and heat flow across a finite temperature difference are irreversible: they create
entropy. A slow, quasistatic compression, by contrast, is reversible: the gas is always in
equilibrium, and any entropy that enters or leaves the gas does so in a way that can be exactly
undone by reversing the process.

% ---- Page 48: Problems 2.28–2.30 ----

\subsection*{Problems}

\textbf{2.28.} Each card in a deck is unique, so there are $52!$ distinct arrangements. The
entropy is therefore
\begin{align*}
  S &= k\ln(52!) \approx k\!\left(52\ln 52 - 52 + \tfrac{1}{2}\ln(2\pi \cdot 52)\right) \\
    &= k\!\left(52\ln 52 - 52 + \ln\!\sqrt{2052}\right) \\
    &\approx 1{,}300\,k \;\approx\; 1.9 \times 10^{-20}\,\text{J/K}.
\end{align*}

\textbf{2.29.} Two Einstein solids interact: solid $A$ has $N_A = 100$ oscillators and solid $B$
has $N_B = 200$ oscillators, with $q = 100$ total energy units. The total multiplicity for a given
value of $q_A$ is
\[
  \Omega(q_A) = \binom{q_A + 99}{q_A}\binom{(100 - q_A) + 199}{100 - q_A}.
\]

The most likely macrostate ($q_A = 60$, $q_B = 40$, proportional to oscillator numbers):
\[
  \Omega_c = \binom{159}{60}\binom{239}{40} = \frac{159!\;239!}{60!\;99!\;40!\;199!}
  \;\Rightarrow\; S_c \approx 264.4\,k.
\]

At equal sharing ($q_A = 50$, $q_B = 50$):
\[
  \Omega_2 = \binom{149}{50}\binom{249}{50} = \frac{149!\;249!}{50!\;99!\;50!\;199!}
  \;\Rightarrow\; S_2 \approx 263.2\,k.
\]

The least likely macrostate (all energy in $B$, $q_A = 0$):
\[
  S_{\min} = k\ln\binom{299}{100} = k\ln\frac{299!}{100!\;199!} \approx 187.5\,k.
\]

Over long time scales the system explores all macrostates, but spends the vast majority of its
time near the peak. The entropy associated with the full accessible range is
$\sim k\ln\dbinom{399}{100} \approx 262\,k$.

\textbf{Problem 2.30.}
\begin{enumerate}
  \item[a)] Using the result of Problem~2.22, the entropy of a two-state paramagnet at the
  multiplicity peak is
  \[
    S = k\ln\!\frac{2^N}{\sqrt{4\pi N}} \approx 2.77\times10^{23}\,k.
  \]

  \item[b)] The most likely macrostate has multiplicity
  \[
    \Omega_{\text{most likely}} = \frac{(2^N)(2^N)}{\sqrt{4\pi N}} = \frac{2^{2N}}{\sqrt{4\pi N}},
  \]
  so $S_{\text{most likely}} \sim 2.77\times10^{23}\,k$ as well. The timescale required to
  visit every microstate is therefore irrelevant: the entropy computed from the peak multiplicity
  is already essentially exact.
\end{enumerate}

% ---- Pages 49–50: Problem 2.49, Problems 2.33–2.37 ----

\textbf{2.49.} Derivation of the Sackur--Tetrode equation. Starting from
\[
  \Omega_N \approx \frac{1}{N!}\cdot\frac{V^N}{h^{3N}}\cdot
    \frac{\pi^{3N/2}}{\left(\tfrac{3N}{2}\right)!}\,(2mU)^{3N/2},
\]
apply Stirling's approximation to $\ln N! \approx N\ln N - N$ and
$\ln\!\left(\tfrac{3N}{2}\right)! \approx \tfrac{3N}{2}\ln\tfrac{3N}{2} - \tfrac{3N}{2}$:
\begin{align*}
  \Omega_N
  &= \frac{1}{N^N e^{-N}\!\sqrt{2\pi N}}\cdot\frac{V^N}{h^{3N}}
     \cdot\frac{\pi^{3N/2}}{\!\left(\tfrac{3N}{2}\right)^{3N/2}\!e^{-3N/2}\!\sqrt{\tfrac{3N}{2}}}
     \cdot(2mU)^{3N/2} \\[4pt]
  &= \frac{e^N}{N^N}\cdot\frac{V^N}{h^{3N}}
     \cdot\left(\frac{\pi\cdot 2mU}{\tfrac{3N}{2}}\right)^{3N/2}\cdot e^{3N/2} \\[4pt]
  &= \frac{V^N}{N^N h^{3N}}\left(\frac{4\pi mU}{3N}\right)^{3N/2} e^{5N/2} \\[4pt]
  &= \left[\frac{V}{N}\cdot\frac{1}{h^3}\left(\frac{4\pi mU}{3N}\right)^{3/2} e^{5/2}\right]^N \\[4pt]
  &= \left[\frac{V}{N}\left(\frac{4\pi mU}{3Nh^2}\right)^{3/2} e^{5/2}\right]^N. \checkmark
\end{align*}
Taking the logarithm recovers the Sackur--Tetrode equation.

\textbf{2.33.}
\begin{enumerate}
  \item[a)] Argon is heavier than helium ($m_{\text{Ar}} \gg m_{\text{He}}$), so the momentum
  hypersphere of radius $\sqrt{2mU}$ is larger for argon at the same energy. A larger momentum
  surface area means more accessible microstates and hence a greater multiplicity and entropy.
\end{enumerate}

\textbf{2.34.} In an isothermal expansion the internal energy $U = \tfrac{3}{2}NkT$ is fixed
(since $T$ is constant), so the Sackur--Tetrode equation gives
\[
  \Delta S = Nk\ln\frac{V_f}{V_i}.
\]
The heat absorbed in a quasistatic isothermal expansion equals the work done by the gas:
$Q = NkT\ln(V_f/V_i)$, so
\[
  \Delta S = \frac{Q}{T}.
\]
In a free expansion, by contrast, $Q = 0$ yet $\Delta S = Nk\ln(V_f/V_i) > 0$. The formula
$\Delta S = Q/T$ does not apply to irreversible processes; the entropy increase in free expansion
is genuine but has no corresponding heat flow.

\textbf{2.35.} The Sackur--Tetrode entropy becomes negative when the argument of the logarithm
falls below $e^{-5/2}$, i.e.\ when
\[
  \frac{V}{N}\left(\frac{4\pi mU}{3Nh^2}\right)^{3/2} < e^{-5/2}.
\]
Solving for $U$:
\[
  \frac{4\pi mU}{3Nh^2} = \left(\frac{N}{V}\right)^{2/3} e^{-5/3}
  \;\Longrightarrow\;
  U = \frac{3Nh^2}{4\pi m}\left(\frac{N}{V}\right)^{2/3} e^{-5/3}.
\]
A negative value of $S$ signals that the Sackur--Tetrode equation has broken down: the gas is
so cold and dense that quantum effects dominate and the classical counting of microstates is no
longer valid.

\textbf{2.37.} A mixture of two gases $A$ and $B$ contains $(1-x)N$ molecules of $A$ and $xN$
molecules of $B$. Assuming both species have the same mass and energy per molecule, and that the
total volume is $V_{\text{tot}}$, the partial volumes before mixing are
\[
  V_A = (1-x)\,V_{\text{tot}}, \qquad V_B = x\,V_{\text{tot}}.
\]
After mixing, each species occupies the full volume $V_{\text{tot}}$, and the entropy of mixing is
\[
  \Delta S_{\text{mix}} = -Nk\bigl[(1-x)\ln(1-x) + x\ln x\bigr],
\]
which is non-negative for all $0 \leq x \leq 1$ and vanishes only for pure species ($x = 0$ or
$x = 1$).
